\documentclass[12pt,a4paper]{article}

% ── Packages ─────────────────────────────────────────────────────
\usepackage[utf8]{inputenc}
\usepackage[T1]{fontenc}
\usepackage{amsmath,amssymb}
\usepackage{graphicx}
\usepackage[margin=2.5cm]{geometry}
\usepackage{hyperref}
\usepackage{booktabs}
\usepackage{caption}
\usepackage{subcaption}
\usepackage{float}
\usepackage{natbib}
\usepackage{enumitem}
\usepackage{xcolor}
\usepackage{listings}
\usepackage{fancyhdr}

% ── Code listing style ───────────────────────────────────────────
\lstset{
  language=Python,
  basicstyle=\ttfamily\small,
  keywordstyle=\color{blue!70!black},
  commentstyle=\color{green!50!black},
  stringstyle=\color{red!60!black},
  frame=single,
  breaklines=true,
  numbers=left,
  numberstyle=\tiny\color{gray},
  backgroundcolor=\color{gray!5},
}

% ── Header / Footer ─────────────────────────────────────────────
\pagestyle{fancy}
\fancyhf{}
\rhead{Complexity Science Project}
\lhead{Power Grid Cascades \& SOC}
\cfoot{\thepage}

\begin{document}

% ── Custom Title Page ───────────────────────────────────────────
\begin{titlepage}
    \centering
    \vspace*{1cm}
    
    % Logo inclusion
    % Adjust the width=0.3\textwidth parameter to scale your logo properly
    \includegraphics[width=0.3\textwidth]{Indian_Institute_of_Technology_Delhi_Logo.svg.png}\\[1cm]
    
    % Institution Information
    {\scshape\Large Indian Institute of Technology Delhi \par}
    \vspace{0.4cm}
    {\large Department of Chemical Engineering \par}
    
    \vspace{2cm}
    
    % Title Section with Horizontal Rules
    \rule{\textwidth}{1pt}\\[0.5cm]
    {\huge \bfseries Network Dynamics and Self-Organized Criticality\\[0.3em] 
    in Power Grid Failure Cascades \par}
    \vspace{0.4cm}
    \rule{\textwidth}{1pt}\\[1.5cm]
    
    % Subtitle / Course Info
    {\Large \textit{CLL798: Complexity Science Individual Project} \par}
    
    \vspace{2.5cm}
    
    % Author Information
    {\Large Md Tahseen Ahmad}\\[0.3em]
    {\large 2022CH11461}\\[0.3em]
    \texttt{ch1221461@iitd.ac.in}
    
    \vfill
    
    % Date
    {\large April 2026 \par}
\end{titlepage}

% ═════════════════════════════════════════════════════════════════
\begin{abstract}
Power grids are among the most critical infrastructure systems in modern society, and their failure dynamics exhibit hallmarks of complex systems.
In this work, we model a power transmission network as a graph of interconnected substations and analyse how small, random perturbations can trigger cascading blackouts that propagate through the network.
Using a sandpile-inspired drive-dissipation model on scale-free, small-world, and random network topologies, we simulate 20,000 time steps of stochastic load injection and cascade relaxation.
The resulting avalanche-size distribution follows a power law with exponent $\alpha \approx 2.31$ on the scale-free topology, consistent with self-organized criticality (SOC).
We further examine the relationship between network connectivity and cascade severity, finding that large avalanches coincide with sharp drops in the size of the largest connected component.
Comparative analysis across topologies reveals that scale-free networks are most susceptible to heavy-tailed cascade distributions, while random networks exhibit thinner tails.
These findings support the hypothesis that real power grids, which share structural features with scale-free networks, naturally evolve towards a self-organized critical state.
\end{abstract}

\textbf{Keywords:} self-organized criticality, power grid, cascading failure, avalanche dynamics, complex networks, sandpile model


% ═════════════════════════════════════════════════════════════════
\section{Introduction}
\label{sec:intro}

Modern electrical power grids are vast, interconnected networks comprising thousands of generators, substations, and transmission lines.
Despite decades of engineering optimization, large-scale blackouts continue to occur with surprising regularity.
The 2003 Northeast blackout in the United States affected 55 million people; the 2012 India blackout left over 600 million without power.
These events share a common pattern: a small, seemingly insignificant triggering event leads to a disproportionately large system-wide failure through a chain reaction of overloads.

This pattern---small perturbations producing scale-invariant cascading failures---is a hallmark of \textit{self-organized criticality} (SOC), a concept introduced by Bak, Tang, and Wiesenfeld \cite{bak1987} through their abelian sandpile model.
In SOC systems, the dynamics of slow driving and fast relaxation push the system towards a critical state without the need for external parameter tuning.
At criticality, the system exhibits power-law distributed event sizes, long-range correlations, and sensitivity to perturbations at all scales.

Power grids are compelling candidates for complexity science analysis for several reasons:

\begin{enumerate}[leftmargin=2em]
  \item \textbf{Noise}: Random fluctuations in demand, renewable generation intermittency, equipment degradation, and weather events continuously perturb the system.
  \item \textbf{Avalanches}: When a transmission line or substation is overloaded, its load is redistributed to neighbouring components, which may themselves become overloaded, triggering a cascade.
  \item \textbf{Connectivity}: The network topology---how substations are linked---determines how perturbations propagate. Hub nodes (high-degree substations) can either stabilize the grid by absorbing load or amplify failures by transmitting overloads to many neighbours simultaneously.
\end{enumerate}

In this project, we develop a computational model that captures these three elements and investigate whether the simulated power grid self-organizes to a critical state.
The remainder of this paper is organized as follows: Section~\ref{sec:background} reviews the theoretical background; Section~\ref{sec:model} describes the model and simulation methodology; Section~\ref{sec:results} presents and analyses the results; Section~\ref{sec:discussion} discusses the implications; and Section~\ref{sec:conclusion} offers conclusions.


% ═════════════════════════════════════════════════════════════════
\section{Theoretical Background}
\label{sec:background}

\subsection{Self-Organized Criticality}

Self-organized criticality (SOC) was proposed by Bak, Tang, and Wiesenfeld in 1987 \cite{bak1987} to explain the ubiquity of power-law behaviour in nature.
The canonical SOC model is the \textit{sandpile}: grains of sand are added one at a time to a pile; when the local slope exceeds a threshold, an avalanche occurs, redistributing grains to neighbours.
The system naturally evolves to a critical state where the avalanche-size distribution follows a power law:
\begin{equation}
  P(s) \sim s^{-\alpha}
  \label{eq:powerlaw}
\end{equation}
where $s$ is the avalanche size and $\alpha$ is the critical exponent, typically in the range $1 < \alpha < 3$ for physical SOC systems.

Three necessary conditions for SOC are:
\begin{enumerate}[leftmargin=2em]
  \item A slow \textit{driving} mechanism that adds energy or stress to the system.
  \item A fast \textit{relaxation} mechanism (avalanches) that dissipates the accumulated stress.
  \item A \textit{threshold} condition that separates stable from unstable states.
\end{enumerate}

\subsection{Complex Networks and Power Grids}

Real power grids exhibit heterogeneous degree distributions that lie between purely random (Erd\H{o}s--R\'enyi) and scale-free (Barab\'asi--Albert) topologies \cite{albert2004}.
Scale-free networks, characterized by a power-law degree distribution $P(k) \sim k^{-\gamma}$, are known to be robust against random failures but highly vulnerable to targeted attacks on high-degree hub nodes \cite{barabasi1999}.

The \textit{connectivity} of the network---measured by the relative size of the largest connected component---is a macroscopic order parameter that reflects system integrity.
A sharp drop in connectivity during a cascade signals a large-scale fragmentation event analogous to a percolation transition.

\subsection{Cascading Failures as Avalanches}

In the power grid context, an avalanche proceeds as follows:
\begin{enumerate}[leftmargin=2em]
  \item A node (substation) becomes overloaded when its load exceeds its capacity.
  \item The overloaded node fails and redistributes a fraction of its load to neighbouring alive nodes.
  \item Neighbours that were already near capacity may themselves become overloaded, extending the cascade.
  \item The cascade halts when no further nodes are overloaded.
\end{enumerate}

This is directly analogous to the sandpile relaxation: the load plays the role of sand grains, the capacity threshold plays the role of the critical slope, and the redistribution to neighbours is the toppling rule.


% ═════════════════════════════════════════════════════════════════
\section{Model and Simulation Methodology}
\label{sec:model}

\subsection{Network Construction}

We model the power grid as an undirected graph $G = (V, E)$ with $N = 200$ nodes and edges generated according to three topologies:

\begin{itemize}[leftmargin=2em]
  \item \textbf{Scale-free} (Barab\'asi--Albert, $m=3$): preferential attachment generates a heterogeneous degree distribution with hub nodes, mimicking real grid structure.
  \item \textbf{Small-world} (Watts--Strogatz, $k=6$, $p=0.3$): high clustering with short path lengths.
  \item \textbf{Random} (Erd\H{o}s--R\'enyi, $p = 2m/N$): homogeneous degree distribution.
\end{itemize}

Each node $i$ is assigned:
\begin{itemize}[leftmargin=2em]
  \item A \textbf{load} $L_i(t)$, initially zero, representing the current electrical load.
  \item A \textbf{capacity} $C_i = \max(k_i, 1)$, where $k_i$ is the node degree. This captures the engineering principle that highly connected substations are built with proportionally higher capacity.
  \item A binary \textbf{status} flag: alive or failed.
\end{itemize}

\subsection{Dynamics}

The simulation follows a drive-dissipation cycle over $T = 20{,}000$ time steps:

\paragraph{Drive (Noise).}
At each step, a random alive node receives a load increment $\Delta L = 0.25$ (the noise magnitude). This represents stochastic demand fluctuations, minor equipment faults, or weather-induced load changes.

\paragraph{Relaxation (Avalanche).}
After the perturbation, the system relaxes through iterated sweeps:
\begin{enumerate}[leftmargin=2em]
  \item For every alive node $i$ with $L_i > C_i$: node $i$ fails, and a fraction $f = 0.8$ of its load is distributed equally among its alive neighbours:
  \begin{equation}
    L_j \leftarrow L_j + \frac{f \cdot L_i}{|\mathcal{N}_i^{\text{alive}}|}, \quad \forall\, j \in \mathcal{N}_i^{\text{alive}}
    \label{eq:redistribution}
  \end{equation}
  where $\mathcal{N}_i^{\text{alive}}$ is the set of alive neighbours of $i$.
  \item Sweeps repeat until no more nodes are overloaded.
\end{enumerate}

The \textbf{avalanche size} $s$ is the total number of nodes that fail during one relaxation episode.

\paragraph{Slow Repair (Dissipation).}
To prevent the system from reaching total collapse and to sustain the SOC dynamics over long simulations, a slow repair mechanism is introduced: if more than $30\%$ of all nodes have failed, $10\%$ of the failed nodes are randomly repaired (reset to alive with zero load). This mimics the real-world process of grid operators restoring service after outages and keeps the system in the vicinity of the critical state.

\subsection{Measurements}

At each time step, we record:
\begin{itemize}[leftmargin=2em]
  \item The avalanche size $s(t)$.
  \item The number of cascade steps.
  \item The relative size of the largest connected component among alive nodes: $\phi(t) = |C_{\max}| / N$.
\end{itemize}

Post-simulation, we compute:
\begin{itemize}[leftmargin=2em]
  \item The probability density function $P(s)$ and complementary cumulative distribution function (CCDF) of avalanche sizes.
  \item A power-law fit $P(s) = a \cdot s^{-\alpha}$ via nonlinear least squares.
  \item The correlation between avalanche size and connectivity drop $\Delta\phi$.
\end{itemize}

\subsection{Implementation}

The simulation is implemented in Python~3 using \texttt{NetworkX} for graph construction and analysis, \texttt{NumPy} for numerical computation, \texttt{SciPy} for curve fitting, and \texttt{Matplotlib} for visualization.
The complete source code is available at the GitHub repository linked in the submission.


% ═════════════════════════════════════════════════════════════════
\section{Results and Analysis}
\label{sec:results}

\subsection{Network Structure}

The scale-free network topology utilized in the primary simulation features nodes coloured by their load-to-capacity ratio, highlighting hub nodes (high degree) as central, well-connected vertices (Figure~\ref{fig:network}). Furthermore, the degree distribution follows the characteristic heavy-tailed form expected of Barab\'asi--Albert networks (Figure~\ref{fig:degree}).

\begin{figure}[H]
  \centering
  \begin{subfigure}[t]{0.48\textwidth}
    \centering
    \includegraphics[width=\textwidth]{figures/network.png}
    \caption{Network visualization. Colour indicates load/capacity ratio (yellow: low, red: high). Black nodes are currently failed.}
    \label{fig:network}
  \end{subfigure}
  \hfill
  \begin{subfigure}[t]{0.48\textwidth}
    \centering
    \includegraphics[width=\textwidth]{figures/degree_distribution.png}
    \caption{Degree distribution of the Barab\'asi--Albert network ($N=200$, $m=3$).}
    \label{fig:degree}
  \end{subfigure}
  \caption{Structure of the simulated power grid network.}
  \label{fig:structure}
\end{figure}

\subsection{Avalanche Time Series}

Plotting the avalanche size as a function of time reveals that the overwhelming majority of perturbations produce no cascade ($s = 0$) or very small cascades ($s = 1$--$2$) (Figure~\ref{fig:timeseries}). However, intermittent large events with $s > 50$ punctuate the background noise, a pattern consistent with the bursty dynamics characteristic of SOC systems. Out of 20,000 time steps, 633 produced non-zero avalanches, reaching a maximum cascade size of 139 nodes (approximately $70\%$ of the network).

\begin{figure}[H]
  \centering
  \includegraphics[width=0.95\textwidth]{figures/avalanche_timeseries.png}
  \caption{Time series of avalanche (cascade) sizes over 20,000 time steps. Note the intermittent large events amid a background of small or zero-size cascades.}
  \label{fig:timeseries}
\end{figure}

\subsection{Avalanche Size Distribution}

The avalanche-size probability distribution, visualized on log--log axes alongside its complementary cumulative distribution function (CCDF), is well approximated by a power law (Equation~\ref{eq:powerlaw}) with an exponent of $\alpha \approx 2.31$, fitted via nonlinear least squares (Figure~\ref{fig:freq}). 

This exponent falls within the range $2 < \alpha < 3$ reported for many SOC systems, including the BTW sandpile model ($\alpha \approx 1.2$--$2.7$ depending on dimensionality) and empirical power grid blackout statistics ($\alpha \approx 1.5$--$2.5$) \cite{carreras2004}. Additionally, the CCDF, which is less sensitive to binning artifacts, confirms the heavy-tailed nature of the distribution.

\begin{figure}[H]
  \centering
  \includegraphics[width=0.95\textwidth]{figures/freq_distribution.png}
  \caption{Left: Avalanche-size probability density on log--log axes with power-law fit ($\alpha = 2.31$). Right: Complementary cumulative distribution function (CCDF).}
  \label{fig:freq}
\end{figure}

\subsection{Connectivity Dynamics}

Tracking the time evolution of the largest connected component fraction, $\phi(t)$, yields a mean connectivity of $\phi \approx 0.735$, indicating that on average about $73.5\%$ of the network remains in a single connected cluster (Figure~\ref{fig:connectivity}). Periodic sharp dips in this connectivity correspond to large cascade events that fragment the network.

\begin{figure}[H]
  \centering
  \includegraphics[width=0.95\textwidth]{figures/connectivity.png}
  \caption{Relative size of the largest connected component over time. The dashed line shows the mean value $\phi = 0.735$.}
  \label{fig:connectivity}
\end{figure}

\subsection{Relationship Between Avalanches and Connectivity}

Examining the scatter plot of connectivity drop $\Delta\phi$ versus avalanche size on log--log axes reveals a clear positive correlation: larger avalanches produce larger drops in network connectivity (Figure~\ref{fig:conn_aval}). This confirms that cascading failures do not merely affect individual nodes but systematically degrade the network's ability to transmit power across long distances.

\begin{figure}[H]
  \centering
  \includegraphics[width=0.55\textwidth]{figures/conn_vs_avalanche.png}
  \caption{Connectivity drop versus avalanche size (log--log). Larger cascades produce proportionally larger fragmentation of the network.}
  \label{fig:conn_aval}
\end{figure}

\subsection{Topology Comparison}

A comparison of the avalanche-size CCDF across three network topologies demonstrates that the scale-free network exhibits the heaviest tail, indicating the highest probability of very large cascades (Figure~\ref{fig:topo}). The small-world network shows intermediate behaviour, while the random network produces the lightest tail. 

This ordering reflects the role of hub nodes: in scale-free networks, the failure of a highly connected hub redistributes load to many neighbours simultaneously, creating a multiplicative amplification effect. Random networks, lacking such hubs, distribute load more uniformly and are therefore less susceptible to extreme cascades.

\begin{figure}[H]
  \centering
  \includegraphics[width=0.6\textwidth]{figures/topology_comparison.png}
  \caption{Comparison of avalanche-size CCDFs across scale-free, small-world, and random network topologies.}
  \label{fig:topo}
\end{figure}

% ═════════════════════════════════════════════════════════════════
\section{Discussion}
\label{sec:discussion}

\subsection{Evidence for Self-Organized Criticality}

Our simulation exhibits the three defining signatures of SOC:

\begin{enumerate}[leftmargin=2em]
  \item \textbf{Power-law distributed avalanches}: The cascade-size distribution follows $P(s) \sim s^{-2.31}$ over more than two decades, indicating scale-invariant dynamics with no characteristic event size.
  \item \textbf{Separation of time scales}: The slow driving (one noise perturbation per step) and fast relaxation (cascades complete within the same step) satisfy the classic SOC separation-of-scales requirement.
  \item \textbf{Self-tuning to criticality}: The system reaches and maintains the critical state without any external tuning of a control parameter. The interplay between noise-driven load accumulation, threshold-triggered cascades, and slow repair naturally steers the system to the boundary between stability and collapse.
\end{enumerate}

\subsection{Role of Connectivity}

The positive correlation between avalanche size and connectivity drop (Figure~\ref{fig:conn_aval}) reveals that network structure is not merely a passive substrate but an active participant in the cascade dynamics.
Hub failure can fragment the network into isolated components, preventing further cascade propagation but also severing power delivery to large regions.

The topology comparison (Figure~\ref{fig:topo}) demonstrates that scale-free networks---which most closely resemble real power grids---are precisely the topology most vulnerable to SOC-driven extreme cascades.
This has important engineering implications: designing grids with more homogeneous connectivity or introducing redundant pathways around hubs could reduce the probability of extreme blackouts.

\subsection{Relation to Real Power Grids}

Empirical studies of blackout data from North America and Europe have consistently found power-law distributions in blackout sizes \cite{carreras2004}, with exponents ranging from approximately 1.5 to 2.5.
Our simulated exponent of $\alpha \approx 2.31$ falls within this empirical range, suggesting that despite the model's simplifications, it captures the essential physics of real cascade dynamics.

The model's simplifications---uniform noise magnitude, deterministic redistribution, and absence of generation-demand balancing---limit direct quantitative comparison with real grids.
However, the qualitative agreement in the avalanche statistics supports the hypothesis that the SOC mechanism is a robust, topology-dependent feature of power grid dynamics.

\subsection{Limitations}

Several limitations should be acknowledged.
The network size ($N=200$) is small compared to real power grids (thousands to tens of thousands of nodes).
The load redistribution rule is simplified; real power flow follows Kirchhoff's laws and depends on line impedances.
The repair mechanism is stylized; real grid restoration follows complex operational procedures.
Future work could address these by incorporating AC/DC power flow equations, realistic demand profiles, and larger network sizes.


% ═════════════════════════════════════════════════════════════════
\section{Conclusion}
\label{sec:conclusion}

We have demonstrated through computational simulation that a model power grid, driven by random load perturbations and subject to threshold-based cascading failures, self-organizes to a critical state.
The avalanche-size distribution follows a power law with exponent $\alpha \approx 2.31$, consistent with both theoretical SOC predictions and empirical blackout statistics.

The key findings are:
\begin{enumerate}[leftmargin=2em]
  \item The power grid model exhibits self-organized criticality without parameter tuning.
  \item Scale-free network topology amplifies extreme cascades compared to random or small-world topologies.
  \item Network connectivity drops are positively correlated with cascade sizes, confirming that large blackouts fundamentally restructure the grid topology.
  \item The system's critical state emerges from the interplay of three elements: noise (random demand fluctuations), avalanches (cascading overloads), and connectivity (network structure governing load redistribution).
\end{enumerate}

These results underscore the importance of complexity science perspectives in understanding and mitigating infrastructure failures.
The self-organized critical nature of power grids suggests that conventional engineering approaches focused on preventing individual component failures may be insufficient; instead, systemic strategies that modify the network's propensity for criticality---such as controlled load shedding, topological restructuring, or adaptive capacity allocation---may be more effective at reducing extreme blackout risk.


% ═════════════════════════════════════════════════════════════════
\section*{Acknowledgements}

This project was completed with the assistance of Claude (Anthropic), which was used for code development, manuscript drafting, and analysis guidance.
All simulation code, analysis, and figures were generated programmatically using Python.
The author acknowledges the use of the following open-source libraries: NetworkX, NumPy, SciPy, and Matplotlib.


% ═════════════════════════════════════════════════════════════════
\bibliographystyle{unsrt}
\begin{thebibliography}{9}

\bibitem{bak1987}
P.~Bak, C.~Tang, and K.~Wiesenfeld,
``Self-organized criticality: An explanation of the 1/f noise,''
\textit{Physical Review Letters}, vol.~59, no.~4, pp.~381--384, 1987.

\bibitem{barabasi1999}
A.-L.~Barab{\'a}si and R.~Albert,
``Emergence of scaling in random networks,''
\textit{Science}, vol.~286, no.~5439, pp.~509--512, 1999.

\bibitem{albert2004}
R.~Albert, I.~Albert, and G.~L.~Nakarado,
``Structural vulnerability of the North American power grid,''
\textit{Physical Review E}, vol.~69, no.~2, p.~025103, 2004.

\bibitem{carreras2004}
B.~A.~Carreras, V.~E.~Lynch, I.~Dobson, and D.~E.~Newman,
``Complex dynamics of blackouts in power transmission systems,''
\textit{Chaos}, vol.~14, no.~3, pp.~643--652, 2004.

\bibitem{watts1998}
D.~J.~Watts and S.~H.~Strogatz,
``Collective dynamics of `small-world' networks,''
\textit{Nature}, vol.~393, no.~6684, pp.~440--442, 1998.

\bibitem{motter2002}
A.~E.~Motter and Y.-C.~Lai,
``Cascade-based attacks on complex networks,''
\textit{Physical Review E}, vol.~66, no.~6, p.~065102, 2002.

\bibitem{crucitti2004}
P.~Crucitti, V.~Latora, and M.~Marchiori,
``Model for cascading failures in complex networks,''
\textit{Physical Review E}, vol.~69, no.~4, p.~045104, 2004.

\bibitem{dobson2007}
I.~Dobson, B.~A.~Carreras, V.~E.~Lynch, and D.~E.~Newman,
``Complex systems analysis of series of blackouts: Cascading failure, critical points, and self-organization,''
\textit{Chaos}, vol.~17, no.~2, p.~026103, 2007.

\end{thebibliography}


% ═════════════════════════════════════════════════════════════════
\appendix
\section{List of Prompts Used}
\label{sec:prompts}

The following prompts were used with Claude (Anthropic) during the development of this project:

\begin{enumerate}
    \item ``I want to build a complexity science simulation modeling power grid blackouts as avalanches, inspired by the BTW sandpile model. Can you help me write a Python script? We need to initialize a network ($N = 200$) where nodes have `load' and `capacity' attributes based on their degree.''
    
    \item ``The initial simulation is crashing too quickly. Let's add a `slow repair' mechanism to the loop to keep it near criticality. Please update our code so that if more than 30\% of the network fails, 10\% of those failed nodes are randomly restored to zero load.''
    
    \item ``The simulation ran successfully, and I have the avalanche size data. Can you write a Python snippet to fit a power-law distribution to this data, and help me draft a brief discussion on how this exponent connects back to our SOC theory?''
    
    \item ``I have written rough notes for the theoretical background and compiled the figures. Please help me refine the academic tone and structure this into a formal LaTeX journal manuscript, ensuring all mathematical notation and graph placements are properly formatted.''
\end{enumerate}

\section{Simulation Code}
\label{app:code}

The complete simulation code (\texttt{power\_grid\_soc.py}) is available at the GitHub repository. Key parameters used in the primary simulation:
\\ Github Repo Link: https://github.com/tahseen252/power-grid-soc-simulation

\begin{table}[H]
  \centering
  \caption{Simulation parameters for the primary run.}
  \label{tab:params}
  \begin{tabular}{lll}
    \toprule
    \textbf{Parameter} & \textbf{Symbol} & \textbf{Value} \\
    \midrule
    Number of nodes     & $N$     & 200 \\
    BA attachment param & $m$     & 3 \\
    Network topology    & ---     & Scale-free (Barab\'asi--Albert) \\
    Total time steps    & $T$     & 20,000 \\
    Noise magnitude     & $\Delta L$ & 0.25 \\
    Redistribution frac & $f$     & 0.80 \\
    Repair threshold    & ---     & 30\% nodes failed \\
    Repair fraction     & ---     & 10\% of failed nodes \\
    \bottomrule
  \end{tabular}
\end{table}

\begin{table}[H]
  \centering
  \caption{Summary statistics from the primary simulation.}
  \label{tab:stats}
  \begin{tabular}{ll}
    \toprule
    \textbf{Metric} & \textbf{Value} \\
    \midrule
    Total time steps            & 20,000 \\
    Non-zero avalanches         & 633 \\
    Maximum avalanche size      & 139 \\
    Mean avalanche size ($s>0$) & 3.50 \\
    Median avalanche size       & 1.0 \\
    Power-law exponent $\alpha$ & 2.31 \\
    Mean connectivity $\phi$    & 0.735 \\
    Minimum connectivity        & 0.005 \\
    \bottomrule
  \end{tabular}
\end{table}


\end{document}
